\section{The Directive That Would Fix This}
\label{sec:ch14-the-directive-that-would-fix-this}

\index{semantic nullability!the shape of the missing feature|(}%
The trade at the end of the last section is not a law of nature. It exists
because GraphQL's type modifier has to answer two questions with one character.
\enquote{Can this value be absent in a correct response} and \enquote{can this
value be missing because something broke} are different questions, and \code{!}
answers both at once.

\index{semantic nullability!what a client actually wants}%
A client wants both answers separately. It wants to know that
\code{Session.title} is always there, so its generated types can say so and no
branch is needed. It also wants a response with a broken service in it to arrive
with a hole rather than an empty page. Those are compatible wishes and nothing
in the language expresses them.

\index{semantic nullability!the working group's own summary}%
That is what the specification's nullability work has been about, and
\code{@semanticNonNull} is the name most of the writing about it uses. The name
is doing less work than it looks like. What it refers to has moved twice since
it was proposed, and none of it is in a published specification.

\subsection{Where the Proposal Actually Is}

\index{semantic nullability!the RFC, and its state}%
The RFC that named the directive is pull request 1065 on the specification
repository, opened by Benjie Gillam on 24 November
2023~\autocite{graphqlrfc1065}. It is still open. It has never been merged, it
has never been closed, and in March 2026 it was marked stale; the project's own
RFC tracker lists it at the first of its stages, strawman.

\index{semantic nullability!where the work went instead}%
Reading only that would give the wrong impression, because a stalled pull
request usually means nobody cared. The opposite happened. The GraphQL
Nullability Working Group archived itself in February 2026, and its closing note
says where the work went~\autocite{nullabilitywg}:

\begin{quote}
Many thanks to everyone who contributed the many discussions and that made it
possible to land on the current \textbf{service capabilities} + \code{onError}
proposal.
\end{quote}

\index{onError@\code{onError}!the proposal that replaced the directive}%
That proposal is pull request 1163, by the same author, and it is a different
shape entirely~\autocite{graphqlrfc1163}. Rather than annotating the schema, it
lets the client say, per request, what it wants done with an error: propagate
it in the usual way, or put a null at the field that failed and report the error
beside it. It sits one stage further along, and where 1065 has not been touched
in months except to be marked stale, 1163 was still being edited two days
before I wrote this.

\index{semantic nullability!nothing is in a published edition}%
Neither is in a published edition of the specification. Searching the September
2025 edition and the current working draft for either name finds nothing.

\subsection{What Hot Chocolate 16 Actually Ships}

\index{semantic nullability!an export-time rewrite}%
The directive exists in Hot Chocolate 16, and it is not a runtime feature. It
is a flag on the schema export:

\begin{minted}[fontsize=\footnotesize]{text}
  --semantic-non-null          Rewrite the exported schema to strip non-null
                               wrappers from output fields and apply the
                               @semanticNonNull directive instead.
\end{minted}

\index{semantic nullability!what the rewrite produces}%
Run it and every output field on the query side loses its exclamation mark and
gains a directive, bare for a single value and carrying a level for a list where
only the items were non-null. The rewritten document declares the directive
itself at the bottom. Two things it leaves alone, both of which the name would
lead you to expect it to touch: the mutation root field keeps its non-null
payload, and argument types are untouched throughout. It is a rewrite of what
the query side returns, not of the schema.

\index{semantic nullability!it changes nothing about execution}%
Nothing about the running service changes. The same requests get the same
responses, propagate the same way, and empty the same lists. The document is
for a client's code generator, so that
generated types can say \enquote{present unless something broke} while the wire
keeps the nullable shape that lets a response degrade.

\index{semantic nullability!the router serves it}%
It composes, which surprised me. \code{wgc} takes a rewritten subgraph document
without complaint and the directive survives into the schema the router serves,
so this is a thing a graph could actually adopt. One detail is worth recording
because it is the third time in this book that two vendors have described the
same directive differently: the document Hot Chocolate writes declares the
directive's argument as a nullable list, and the composed schema declares it as
a non-null one. The composer is substituting a definition of its own rather than
carrying through the one it was given.

\index{ErrorHandlingMode@\code{ErrorHandlingMode}!the runtime half}%
The runtime half is in the same release under a different name. Hot Chocolate
16.6.1 reads an \code{onError} property off the request body, taking
\code{PROPAGATE} or \code{NULL}, and there is a server-side default beside a
switch controlling whether a client may override it. Michael Staib describes
what the second value does on ChilliCream's blog~\autocite{staib2026hc16}:

\begin{quote}
With \code{onError: "NULL"}, Hot Chocolate stops null propagation, returns
\code{null} at the field that failed, and still reports the error.
\end{quote}

\index{ErrorHandlingMode@\code{ErrorHandlingMode}!ahead of the specification}%
That is an implementation of pull request 1163, on the wire, in a shipped
release, while the pull request is unmerged and one stage into a process with
several. It also replaced something: version 15 had a server-wide switch that
turned the same non-propagating behavior on for every request, and the
migration notes say it was removed in favor of the proposal. So the setting
moved from the server to the request, which is the difference between an
operator choosing and a client choosing.

\subsection{Why This Book Ships Neither}

\index{semantic nullability!why the graph does not adopt it}%
Both mechanisms live inside one service. \code{onError} is a property of a
request a client sends to a Hot Chocolate server, and every request in this book
since chapter~\ref{ch:enter-the-router} is sent to a router written in Go. What
the router does with an error is the router's, and it forwards a query rather
than a client's envelope. Making that end-to-end is a question about two
products agreeing, and neither of them has answered it in a released version.

\index{semantic nullability!what a reader should do about it}%
So the honest position is the one this chapter has already taken. Put the
question marks where a promise can be kept, accept the null check on the client,
and watch pull request 1163 rather than the directive, because the directive is
the part of this that has already been superseded once. When the two halves
arrive together, the trade at the end of
section~\ref{sec:ch14-nullability-for-a-graph-you-share} stops being a trade,
and a schema will be able to tell a client that a value is always there and that
this response is missing it.

\begin{onhc14}
\index{semantic nullability!absent on 14}%
There is no such flag. On 14.3.1 the schema export command takes an output path
and a schema name and nothing else, and passing the flag anyway is refused by
name:

\begin{minted}{text}
Unrecognized command or argument '--semantic-non-null'.
\end{minted}

Nothing else in this chapter is a version difference. The propagation rule is
the specification's and behaves the same on both. Three of the five fields this
chapter softens belong to services that branch does not carry at all, and the
two on the Sessions service are not ported to it, because
appendix~\ref{app:hc14} keeps that branch to the code a callout quotes and no
callout quotes them.
\end{onhc14}
\index{semantic nullability!the shape of the missing feature|)}%
